\section{Design and estimation}\label{sec:design}

\subsection{Stratification}
Let $K_i=\sum_{m=1}^M\hat y_{mi}\in\{0,\dots,M\}$ be the number of models that predict the positive class for instance $i$. The proposed design, which we call PATTERN, partitions the pool into strata indexed by $(G_i,K_i)$. The two extreme strata of each group, $K_i=0$ and $K_i=M$, contain the instances on which all models agree; the strata with $0<K_i<M$ contain every instance on which at least one pair of models disagrees, and the value of $K_i$ indicates how evenly the models are split. Strata with fewer than 50 instances are merged with the adjacent stratum (in $K$) of the same group, so that within-stratum variances can be estimated; this gives $H\le 2(M+1)$ strata, of size $N_h$.

The count $K_i$ is a coarse summary of the prediction pattern $(\hat y_{1i},\dots,\hat y_{Mi})$, which has up to $2^M$ values. It is chosen because it is invariant to the ordering of models, it keeps the number of strata linear in $M$, and its extreme values isolate the agreement region exactly. Finer stratifiers are compared in Section~\ref{sec:ablations}.

\subsection{Estimation}
Within each stratum a simple random sample of $n_h$ instances is drawn without replacement and labelled ($\sum_h n_h=L$). The total $T_c$ is estimated by the stratified expansion estimator
\begin{equation}\label{eq:HT}
\hat T_c=\sum_{h=1}^H N_h\,\bar t_{hc},\qquad \bar t_{hc}=\frac{1}{n_h}\sum_{i\in s_h} t_{ic}(Y_i),
\end{equation}
which is unbiased for $T_c$ \citep{cochran1977}. The comparisons are estimated by the plug-in $\hat\theta=f(\hat T)$, a ratio-type estimator. Its variance is estimated by Taylor linearisation \citep[Section~5.5]{sarndal1992}: with the influence values $z_{i}=\nabla f(\hat T)^{\top} t_i(Y_i)$,
\begin{equation}\label{eq:var}
\hat V(\hat\theta_k)=\sum_{h=1}^H N_h^2\Big(1-\frac{n_h}{N_h}\Big)\frac{s^2_{z_k,h}}{n_h},
\end{equation}
where $s^2_{z_k,h}$ is the sample variance of the $k$th influence values in stratum $h$. Wald intervals $\hat\theta_k\pm 1.96\,\hat V(\hat\theta_k)^{1/2}$ are reported. If a sampled denominator total is zero (no sampled positive in a group), the estimate is undefined; this happens with small probability under SRS at small budgets and never occurred under PATTERN in our experiments (Section~\ref{sec:empirical}).

\subsection{Allocation}
The allocation $(n_1,\dots,n_H)$ determines the efficiency of the design. For a single estimand the variance of the linearised estimator is minimised by Neyman allocation, $n_h\propto N_h S_h$, where $S_h$ is the within-stratum standard deviation of the influence values. With many estimands we minimise the weighted sum $\sum_k w_k V_k$ of their linearised variances, which leads (Section~\ref{sec:theory}) to
\begin{equation}\label{eq:alloc}
n_h\ \propto\ N_h\Big(\sum_k w_k S^2_{hk}\Big)^{1/2}.
\end{equation}
We use $w_k=1/S_k^2$, where $S_k^2$ is the population variance of the influence values of estimand $k$; this weights each estimand by the inverse of its variance under SRS, so that every comparison counts equally in relative terms.

The stratum variances depend on the unknown labels. We replace them by their expectations under a proxy model in which $Y_i$ is Bernoulli with probability $\pi_i=M^{-1}\sum_m\hat p_{mi}$, the average predicted probability of the models. Because the influence values are linear in $Y_i$, $z_{ik}=Y_iA_{ik}+(1-Y_i)B_{ik}$ with known $A_{ik},B_{ik}$ evaluated at the proxy totals, and the proxy moments $\mathrm E_\pi z_{ik}=\pi_iA_{ik}+(1-\pi_i)B_{ik}$ and $\mathrm E_\pi z^2_{ik}=\pi_iA^2_{ik}+(1-\pi_i)B^2_{ik}$ give closed-form proxy variances. No label is used. A poor proxy can only reduce efficiency; the estimator remains unbiased for the totals because the inclusion probabilities $n_h/N_h$ are known. Each stratum receives at least ten labels (or all of its instances if it is smaller), and rounding to integers is done greedily so that the budget is met exactly.

\begin{algorithm}[t]
\caption{PATTERN design for comparing $M$ classifiers on group-fairness contrasts}\label{alg:pattern}
\begin{algorithmic}[1]
\Require pool $U$ with predictions $\hat y_{mi}$, probabilities $\hat p_{mi}$, groups $G_i$; budget $L$
\State $K_i\gets\sum_m\hat y_{mi}$; strata $h\gets(G_i,K_i)$; merge strata with $N_h<50$ within group
\State $\pi_i\gets M^{-1}\sum_m\hat p_{mi}$; proxy totals $\tilde T\gets\sum_i\{\pi_iu_i+(1-\pi_i)v_i\}$
\State $A_i\gets\nabla f(\tilde T)^\top u_i$, $B_i\gets\nabla f(\tilde T)^\top v_i$; proxy moments of $z_i$
\State $S^2_{hk}$, $S^2_k\gets$ proxy within-stratum and overall variances
\State $n_h\gets$ allocation \eqref{eq:alloc} with $w_k=1/S^2_k$, minimum 10, capped at $N_h$, $\sum_h n_h=L$
\State draw $n_h$ instances without replacement in each stratum and query their labels
\State $\hat T\gets$ \eqref{eq:HT}; $\hat\theta\gets f(\hat T)$; $\hat V\gets$ \eqref{eq:var}; Wald intervals
\State \Return $\hat\theta$, standard errors, intervals
\end{algorithmic}
\end{algorithm}

\subsection{Reference designs}
We compare PATTERN with three designs that use the same estimators. SRS samples $L$ instances without replacement from the whole pool. GROUP stratifies by the group attribute only and allocates with \eqref{eq:alloc}; it oversamples the minority group when that reduces variance, which is the natural first idea for fairness auditing. ORACLE-ALLOC uses the PATTERN strata with the allocation computed from the true labels instead of the proxy. It is not usable in practice and is not an upper bound on efficiency, since \eqref{eq:alloc} minimises a weighted sum of first-order variances rather than the mean squared error of each estimand; it shows how much is lost by using a proxy. Further baselines from the active-testing and survey-sampling literatures are compared in Section~\ref{sec:baselines}.

\subsection{Computational cost}
The design needs one pass over the pool to compute $K_i$, $\pi_i$ and the proxy moments, which costs $O(NMC)$ operations with $C=O(M)$ totals, and the estimation costs $O(LMC)$. For the pools used here ($N\le 19{,}537$, $M=9$) both steps take well under a second on a laptop.
